% =====================================================================
% Fase 1 — Arquitectura física del rack de producción
% Documento maestro. Compilar con:
%     latexmk -pdf main.tex
% =====================================================================
\documentclass[11pt,a4paper]{article}

% ----- Paquetes base --------------------------------------------------
\usepackage[utf8]{inputenc}
\usepackage[T1]{fontenc}
\usepackage[spanish,es-tabla]{babel}
\usepackage{lmodern}
\usepackage{microtype}

% ----- Layout ---------------------------------------------------------
\usepackage[a4paper,margin=2.3cm]{geometry}
\usepackage{parskip}
\usepackage{enumitem}

% ----- Figuras y tablas ----------------------------------------------
\usepackage{graphicx}
\usepackage{float}
\usepackage{booktabs}
\usepackage{array}
\usepackage{longtable}
\usepackage{tabularx}
\usepackage{caption}
\usepackage{subcaption}

% ----- Diagramas vectoriales (TikZ) -----------------------------------
\usepackage{tikz}
\usetikzlibrary{shapes.geometric,arrows.meta,positioning,fit,calc}

% ----- Matemáticas y unidades ----------------------------------------
\usepackage{amsmath,amssymb}
\usepackage{siunitx}
\sisetup{output-decimal-marker={,}}

% ----- Bibliografía ---------------------------------------------------
\usepackage[backend=biber,style=numeric,sorting=none]{biblatex}
\addbibresource{referencias.bib}

% ----- Hyperref (al final) -------------------------------------------
\usepackage[hidelinks]{hyperref}

% ----- Macro para placeholders de cotización -------------------------
% Uso: \cotizar{Descripción exacta del componente y link sugerido}
\newcommand{\cotizar}[1]{%
  \textbf{[COTIZAR:}~\emph{#1}\textbf{]}%
}

% ----- Macro para marcar pendientes post-curso -----------------------
\newcommand{\pendientecurso}[1]{%
  \textbf{[POST-CURSO:}~\emph{#1}\textbf{]}%
}

% ----- Parámetros del diseño -----------------------------------------
% =====================================================================
% params.tex  --  Parámetros editables del diseño físico (Fase 1)
% =====================================================================
% Modifica únicamente este archivo para adaptar el documento a tus
% restricciones reales de espacio, temperatura ambiente y número de
% contenedores. El resto de las secciones consume estos valores.
% =====================================================================

% ---------- Dimensiones disponibles del espacio de instalación -------
\newcommand{\paramHmax}{180}   % altura máxima disponible [cm]
\newcommand{\paramWmax}{90}    % ancho máximo disponible [cm]
\newcommand{\paramDmax}{45}    % fondo máximo disponible [cm]

% ---------- Condiciones ambientales del recinto ----------------------
\newcommand{\paramTambMin}{16} % temperatura ambiente mínima esperada [°C]
\newcommand{\paramTambMax}{26} % temperatura ambiente máxima esperada [°C]
\newcommand{\paramHRamb}{45}   % humedad relativa ambiente típica [%]

% ---------- Configuración del rack -----------------------------------
\newcommand{\paramNiveles}{4}          % número de niveles del rack
\newcommand{\paramNcontenedores}{3}    % número de cámaras de fructificación
\newcommand{\paramVolTote}{55}         % volumen nominal de cada tote [L]
\newcommand{\paramCargaUtilNivel}{50}  % carga útil por nivel mínima recomendada [kg]

% ---------- Especies asignadas (por contenedor) ----------------------
% Asignación tentativa --- VALIDAR POST-CURSO. Ver pendientes-postcurso.md
\newcommand{\paramEspecieA}{Pleurotus ostreatus (oyster)}
\newcommand{\paramEspecieB}{Hericium erinaceus (lion's mane)}
\newcommand{\paramEspecieC}{Lentinula edodes (shiitake)}

% ---------- Sensores por contenedor ----------------------------------
\newcommand{\paramSensorTH}{SHT35}
\newcommand{\paramSensorCO}{MH-Z19B}
\newcommand{\paramSensorLux}{BH1750}
\newcommand{\paramSensorSubstrato}{DS18B20 (opcional)}

% ---------- Esterilización e inoculación (§1.7) ----------------------
\newcommand{\paramVolOllaPresion}{23}    % capacidad olla de presión [L]
\newcommand{\paramPresionEster}{15}      % presión de esterilización [psi]
\newcommand{\paramTiempoEsterSustrato}{90} % tiempo esterilización sustrato [min]
\newcommand{\paramTiempoEsterGrano}{60}    % tiempo esterilización grano [min]

% ---------- Identificador de revisión --------------------------------
\newcommand{\paramRevision}{0.2-relleno-parcial}
\newcommand{\paramFecha}{\today}


% =====================================================================
\title{Sistema integral de cultivo gourmet de hongos\\
       \large Fase 1 --- Arquitectura física del rack de producción}
\author{Hola Tofu --- I+D}
\date{Revisión \paramRevision\ --- \paramFecha}

\begin{document}
\maketitle
\tableofcontents
\clearpage

% ---------------------------------------------------------------------
% =====================================================================
% §1.1 — Plano dimensionado del rack
% =====================================================================
\section{Plano dimensionado del rack}
\label{sec:rack}

Esta sección define la geometría del rack metálico que aloja los
\paramNcontenedores{} contenedores de fructificación. Las dimensiones
máximas disponibles (\paramHmax\,cm $\times$ \paramWmax\,cm $\times$
\paramDmax\,cm) y la temperatura ambiente esperada
(\paramTambMin--\paramTambMax\,\si{\celsius}) son paramétricas y viven
en \texttt{params.tex}.

% ---------------------------------------------------------------------
\subsection{Cálculo de carga útil}
\label{sec:rack:carga}

La carga por nivel se estima como:

\begin{equation}
  m_{\text{nivel}} = m_{\text{tote}} + m_{\text{sustrato}} + m_{\text{accesorios}}
\end{equation}

\noindent con:
\begin{itemize}[leftmargin=*,itemsep=0pt]
  \item $m_{\text{tote}} \approx \SI{2}{\kilogram}$ (HDPE \paramVolTote\,L vacío).
  \item $m_{\text{sustrato}} \approx \SI{12}{\kilogram}$ a \SI{15}{\kilogram}
        (paja/pellets hidratados con HR $\approx 65$\,\%).
  \item $m_{\text{accesorios}} \approx \SI{2}{\kilogram}$ (humidificador,
        ventilador, cableado, sensores).
\end{itemize}

Total estimado por nivel cargado: $\approx$~\SI{17}{\kilogram}. Se
especifica un rack con capacidad declarada
\textbf{$\geq$~\paramCargaUtilNivel{}\,kg por nivel} para conservar un
factor de seguridad $\approx 3$ frente a deformación.

\subsection{Geometría general}
\label{sec:rack:geometria}

Asumiendo \paramNiveles{} niveles distribuidos uniformemente sobre
\paramHmax\,cm de alto:

\begin{equation}
  h_{\text{libre}} = \frac{\text{\paramHmax}~\text{cm}}{\text{\paramNiveles}} \approx 45~\text{cm}
\end{equation}

Un tote estándar de \paramVolTote\,L mide aproximadamente
$60 \times 40 \times 35$\,cm, por lo que la altura libre por entrepaño
debe ser al menos $35 + 10 = 45$\,cm para permitir extracción vertical
sin desensamblar el rack. Si el operador requiere $\geq 1$ nivel de
servicio (electrónica, broker MQTT, humidificadores compartidos), se
reserva el nivel inferior y se reduce a $N - 1$ niveles productivos.

\paragraph{Distribución vertical recomendada (revisión 0.2).}
\begin{itemize}[leftmargin=*,itemsep=0pt]
  \item Nivel 1 (superior, naturalmente más cálido): \paramEspecieB{}
        --- \emph{requiere validación post-curso}.
  \item Nivel 2: \paramEspecieA{}.
  \item Nivel 3: \paramEspecieC{} (preferencia térmica más baja).
  \item Nivel 4 (inferior): nivel de servicio (electrónica, fuentes,
        depósito de agua del humidificador, broker MQTT).
\end{itemize}

% ---------------------------------------------------------------------
\subsection{Selección de materiales}
\label{sec:rack:materiales}

\begin{table}[H]
  \centering
  \caption{Comparativa de materiales para la estructura del rack.}
  \label{tab:rack:materiales}
  \begin{tabularx}{\linewidth}{lXXX}
    \toprule
    \textbf{Material} & \textbf{Pros} & \textbf{Contras} & \textbf{Veredicto} \\
    \midrule
    Acero galvanizado pintado &
      Bajo costo, alta disponibilidad MX (Home Depot). &
      Corrosión local si se raspa la pintura en ambiente HR$>$85\,\%. &
      \textbf{Opción económica.} \\
    Acero inoxidable AISI~304 &
      Resistencia a HR alta y sanitización con peróxido. &
      $\sim 3$--$5\times$ el costo del galvanizado. &
      \textbf{Opción profesional.} \\
    Estantería plástica industrial (PP/HDPE) &
      Inmune a corrosión, ligera. &
      Capacidad de carga típicamente $<$\,30\,kg/nivel; flexión
      bajo peso húmedo. &
      Descartada para esta carga. \\
    \bottomrule
  \end{tabularx}
\end{table}

\subsection{Criterios de aceptación}
\label{sec:rack:aceptacion}

\begin{itemize}[leftmargin=*,itemsep=0pt]
  \item Carga declarada del rack $\geq$~\paramCargaUtilNivel{}\,kg por
        nivel.
  \item Altura libre por entrepaño $\geq$~\SI{45}{\centi\meter}.
  \item Anclaje antivuelco a muro instalado (mandatorio si el rack
        descansa sobre superficie sin nivelar).
  \item Nivel inferior reservado para servicios.
\end{itemize}

% =====================================================================
% §1.2 — Diseño de los contenedores (cámaras de fructificación)
% =====================================================================
\section{Diseño de los contenedores de fructificación}
\label{sec:contenedores}

Cada uno de los \paramNcontenedores{} contenedores herméticos
(\paramVolTote\,L, HDPE alimenticio) se opera como una cámara
independiente tipo \emph{Monotub} modificada, con intercambio gaseoso
pasivo regulado por perforaciones cubiertas de membrana microporosa.

% ---------------------------------------------------------------------
\subsection{Patrón de perforación para intercambio gaseoso pasivo}
\label{sec:contenedores:perforacion}

El patrón se parametriza por especie. La demanda metabólica de
O$_2$ y la tolerancia a CO$_2$ definen el porcentaje de área abierta
sobre la pared lateral del tote~\cite{stamets2000}.

\begin{table}[H]
  \centering
  \caption{Patrón de perforación por especie. Valores de partida
    sujetos a validación experimental con sensor MH-Z19B (Fase 2).}
  \label{tab:perforacion}
  \begin{tabularx}{\linewidth}{Xcccc}
    \toprule
    \textbf{Especie} & $\varnothing$ [mm] & Sep. [mm] & N. orif. & \% área abierta \\
    \midrule
    \paramEspecieA & 6 & 50 & 80--120 & $\approx 3$\,\% \\
    \paramEspecieB & 6 & 70 & 40--60  & $\approx 1{,}5$\,\% \\
    \paramEspecieC & 6 & 70 & 40--60  & $\approx 1{,}5$\,\% \\
    \bottomrule
  \end{tabularx}
\end{table}

Cada perforación se cubre por el interior con cinta de polímero
microporoso (Tyvek 1073B o equivalente) para mitigar el ingreso de
esporas contaminantes mientras se preserva el intercambio gaseoso. El
patrón se distribuye sobre la \textbf{mitad superior} de la pared
lateral, evitando la zona en contacto con sustrato húmedo.

\subsection{Pasamuros para instrumentación y servicios}
\label{sec:contenedores:pasamuros}

Cada contenedor requiere los pasamuros listados en la
tabla~\ref{tab:pasamuros}, sellados con prensaestopas roscados y
arandela EPDM, complementados con cordón de silicona neutra de grado
alimenticio.

\begin{table}[H]
  \centering
  \caption{Pasamuros por contenedor.}
  \label{tab:pasamuros}
  \begin{tabularx}{\linewidth}{lXll}
    \toprule
    \textbf{Servicio} & \textbf{Ubicación} & \textbf{$\varnothing$ cable} & \textbf{Prensaestopa} \\
    \midrule
    Sensor SHT35           & Pared lateral, $\sfrac{2}{3}$ altura      & 4 hilos JST       & PG7 \\
    Sensor MH-Z19B (UART)  & Pared opuesta a SHT35, mismo plano        & 4 hilos UART      & PG9 \\
    Sensor BH1750 (I$^2$C) & Tapa interior, sombra del LED evitada     & 4 hilos JST       & PG7 \\
    Cámara ESP32-CAM       & Tapa superior, centrada                   & flat cable + USB  & PG11 \\
    Manguera humidificador & Pared corta, $\sfrac{1}{4}$ inferior      & manguera 6\,mm OD & PG11 \\
    DS18B20 (sustrato)     & Pared corta opuesta, atravesando sustrato & 3 hilos blindado  & PG7 \\
    \bottomrule
  \end{tabularx}
\end{table}

Total: 6 pasamuros por contenedor $\times$ \paramNcontenedores{}
contenedores $=$ 18 pasamuros consolidados en la BOM (\S\ref{sec:bom}).

\subsection{Sellado e higiene del microambiente}
\label{sec:contenedores:sellado}

\begin{itemize}[leftmargin=*,itemsep=0pt]
  \item Empaque perimetral de silicona alimenticia grado FDA, aplicado
        en frío en el canal de cierre de la tapa.
  \item Sanitización entre lotes: peróxido de hidrógeno al 3\,\% en
        spray, seguido de alcohol isopropílico al 70\,\%. Tiempo de
        aireación mínimo: 30\,min con tapa abierta antes de
        reintroducir bolsa colonizada.
  \item \textbf{Prohibido cloro / hipoclorito:} corrosión inducida sobre
        anclajes del rack y oxidación de prensaestopas.
\end{itemize}

\subsection{Criterios de aceptación}
\label{sec:contenedores:aceptacion}

\begin{itemize}[leftmargin=*,itemsep=0pt]
  \item Cada contenedor mantiene HR $\geq$\,85\,\% durante al menos 4\,h
        sin actuación del humidificador (verificable con telemetría
        de Fase 2).
  \item Los pasamuros no presentan fuga visible al ensayo de presión
        positiva ligera.
  \item La tapa abre y cierra sin re-aplicar silicona durante al menos
        30 ciclos antes de mantenimiento.
\end{itemize}

% =====================================================================
% §1.3 — Asignación de especies por contenedor
% =====================================================================
\section{Asignación de especies por contenedor}
\label{sec:especies}

\textbf{Nota de revisión.} La selección final de especies y sus
\emph{setpoints} óptimos están sujetos a validación tras el curso de
cultivo del operador. Esta sección consolida valores de partida
documentados en literatura para tres especies de uso comercial común.
La tabla \ref{tab:especies} sirve como base para los \emph{setpoints}
del firmware (Fase~2).

% ---------------------------------------------------------------------
\subsection{Parámetros ambientales por especie (literatura)}
\label{sec:especies:tabla}

\begin{table}[H]
  \centering
  \caption{Rangos ambientales óptimos de fructificación reportados en
    literatura. Adaptado de Stamets~\cite{stamets2000} y Chang \&
    Miles~\cite{chang2004}.}
  \label{tab:especies}
  \begin{tabularx}{\linewidth}{Xccccc}
    \toprule
    \textbf{Especie} & T [\si{\celsius}] & HR [\%] & CO$_2$ [ppm] & Lux & Fotop. [h] \\
    \midrule
    \paramEspecieA & 18--24 & 85--95 & $<$1000 & 500--1500 & 12 \\
    \paramEspecieB & 18--24 & 90--95 & $<$800  & 500--1000 & 12 \\
    \paramEspecieC & 12--18 & 80--90 & $<$1000 & 500--2000 & 12 \\
    \bottomrule
  \end{tabularx}
\end{table}

\paragraph{Notas operativas.}
\begin{itemize}[leftmargin=*,itemsep=0pt]
  \item \textbf{Oyster} (\paramEspecieA): especie más permisiva y rápida.
        Ideal como especie de entrenamiento operativo. Tolera CO$_2$
        más alto a costa de elongación del estípite.
  \item \textbf{Lion's mane} (\paramEspecieB): sensible a corrientes
        directas de aire. Forma "barbas" con HR sostenida $>$\,90\,\%;
        bajo deshidratación produce ramificación corta.
  \item \textbf{Shiitake} (\paramEspecieC): prefiere temperaturas más
        bajas y choque térmico inicial (cold shock $\approx 10$\,\si{\celsius}
        durante 24--48\,h) para inducir pinning.
\end{itemize}

\subsection{Asignación propuesta (tentativa)}
\label{sec:especies:asignacion}

Justificación de la asignación contenedor$\leftrightarrow$especie
basada en la jerarquía térmica natural del rack (calor asciende):

\begin{table}[H]
  \centering
  \caption{Asignación propuesta de especies por nivel.}
  \label{tab:asignacion}
  \begin{tabularx}{\linewidth}{cXl}
    \toprule
    \textbf{Nivel} & \textbf{Especie} & \textbf{Justificación} \\
    \midrule
    1 (sup.) & \paramEspecieB & T más alta naturalmente; banda 18--24\,\si{\celsius}. \\
    2        & \paramEspecieA & Banda 18--24\,\si{\celsius}; alta rotación, acceso medio. \\
    3        & \paramEspecieC & Banda térmica más baja del rack (12--18\,\si{\celsius}). \\
    4 (inf.) & --- & Servicios (electrónica, fuentes, broker MQTT). \\
    \bottomrule
  \end{tabularx}
\end{table}

\pendientecurso{Validar la jerarquía térmica del rack midiendo en sitio
con termómetro IR antes de fijar la asignación. Reasignar si la
estratificación real difiere de la asumida.}

\subsection{Riesgo de contaminación cruzada}
\label{sec:especies:contaminacion}

La segregación física por contenedor mitiga contaminación cruzada
entre especies; sin embargo, durante el manejo (apertura de tapa,
cosecha) se generan corrientes que pueden transportar esporas. Medidas
operativas:

\begin{itemize}[leftmargin=*,itemsep=0pt]
  \item Abrir un contenedor a la vez.
  \item Manejar primero la especie con mayor sensibilidad
        (lion's mane) y al final la más robusta (oyster).
  \item Lavarse manos y desinfectar superficie de trabajo entre
        manipulaciones.
\end{itemize}

\subsection{Criterios de aceptación}
\label{sec:especies:aceptacion}

\begin{itemize}[leftmargin=*,itemsep=0pt]
  \item Tabla de parámetros respaldada por al menos dos referencias
        bibliográficas verificables.
  \item Asignación contenedor$\leftrightarrow$especie revalidada con
        mapeo térmico en sitio antes de cargar el primer lote.
  \item Setpoints de Fase~2 derivan directamente de
        tabla~\ref{tab:especies}.
\end{itemize}

% =====================================================================
% §1.4 — Lista de materiales (BOM) — Fase 1 (estructura física)
% =====================================================================
\section{Lista de materiales (BOM) --- Fase 1}
\label{sec:bom}

\subsection*{Resumen de la sección}
Solo se incluyen los componentes \textbf{físicos del rack y los
contenedores}. El BOM electrónico (sensores, ESP32, cámaras,
humidificadores) se consolida en la Fase 8 después de cerrar las
decisiones de las Fases 2 y 3, para evitar re-cotización.

Los precios están deliberadamente vacíos: cada renglón se marca con
\cotizar{descripción exacta a buscar} y se replica en
\texttt{cotizaciones\_pendientes.md} con la búsqueda sugerida por
proveedor (MercadoLibre MX, Home Depot, Steren, AliExpress).

% ---------------------------------------------------------------------
\subsection{Estructura del rack}
\label{sec:bom:rack}

% TODO[Fase 1 — relleno]:
% Tabla con columnas:
%   | Ítem | Descripción | Cant. | Económico (MXN) | Profesional (MXN) | Proveedor sugerido |
% Renglones esperados:
%   - Rack metálico \paramNiveles{} niveles, \paramHmax x \paramWmax x \paramDmax cm
%   - Anclaje a pared (opcional, antivuelco)
%   - Niveladores de patas (4 uds)
%   - Tapete antiderrame para nivel inferior
% Usar la macro \cotizar{...} en lugar de inventar precios.

\subsection{Contenedores y sellado}
\label{sec:bom:contenedores}

% TODO[Fase 1 — relleno]:
% Renglones esperados:
%   - Tote hermético \paramVolTote\,L con tapa (x\paramNcontenedores)
%   - Empaque de silicona alimenticia grado FDA (rollo)
%   - Prensaestopas PG7/PG9/PG11 según pasamuros
%   - Arandelas EPDM
%   - Tornillería inoxidable M3/M4 para fijación de sensores
%   - Tyvek o filtro microporoso para perforaciones (m^2)

\subsection{Herramientas específicas requeridas}
\label{sec:bom:herramientas}

% TODO[Fase 1 — relleno]:
%   - Sierra cilíndrica (hole saw) de diámetros derivados de §1.2.1
%   - Punta cónica para abrir pasamuros en plástico HDPE
%   - Termómetro infrarrojo para mapeo térmico del rack
%   - Higrómetro de referencia para calibración cruzada con SHT35
% Marcar las que ya tengo (laboratorio I+D) vs. las que faltan comprar.

\subsection{Política de placeholders de precio}
\label{sec:bom:politica}

\textbf{Regla estricta:} ningún precio se inventa.
\begin{itemize}[leftmargin=*]
  \item Si el componente fue cotizado por el operador, se sustituye el
        \cotizar{} por el valor real con fecha y URL.
  \item Mientras no haya cotización, el renglón permanece marcado y
        aparece replicado en \texttt{cotizaciones\_pendientes.md}.
  \item Las dos columnas (económico / profesional) se consolidan solo
        cuando ambas tienen al menos una cotización real cada una.
\end{itemize}

\subsection{Criterios de aceptación}
\label{sec:bom:aceptacion}

\begin{itemize}[leftmargin=*]
  \item Todo renglón sin precio aparece en
        \texttt{cotizaciones\_pendientes.md} con búsqueda sugerida.
  \item La cantidad de cada ítem está justificada (no se piden
        "varios", se piden números concretos derivados de \S\ref{sec:rack}
        y \S\ref{sec:contenedores}).
\end{itemize}

% =====================================================================
% §1.5 — Diagrama de flujo del proceso productivo
% =====================================================================
\section{Flujo del proceso productivo}
\label{sec:flujo}

El flujo abarca desde la preparación e inoculación del sustrato hasta
la distribución del producto final, fresco o deshidratado.

% ---------------------------------------------------------------------
\subsection{Etapas y duración orientativa}
\label{sec:flujo:etapas}

\begin{table}[H]
  \centering
  \caption{Etapas del proceso productivo. Tiempos orientativos para
    \paramEspecieA{}, ajustables por especie (\S\ref{sec:especies}).}
  \label{tab:flujo}
  \begin{tabularx}{\linewidth}{lXl}
    \toprule
    \textbf{Etapa} & \textbf{Acción} & \textbf{Duración} \\
    \midrule
    1 & Hidratación del sustrato y mezcla con suplementos.            & 12--24\,h \\
    2 & Embolsado en bolsa autoclavable con parche filtro.            & 0{,}5\,h/lote \\
    3 & Esterilización en olla de presión (\S\ref{sec:preparacion}).  & \paramTiempoEsterSustrato\,min + enfriado 12\,h \\
    4 & Inoculación con grain spawn en SAB.                           & 0{,}3\,h/bolsa \\
    5 & Incubación en oscuridad, T ambiente.                          & 10--21\,días \\
    6 & Inducción de fructificación (al contenedor sellado).          & 0\,h (transición) \\
    7 & Pinning + maduración.                                         & 5--14\,días \\
    8 & Cosecha.                                                      & 0{,}5\,h/contenedor \\
    9a& Empaque fresco $\to$ refrigeración 2--4\,\si{\celsius}.       & inmediato \\
    9b& Deshidratado (\S\ref{sec:postcosecha}) $\to$ empaque hermético.& 6--14\,h \\
    \bottomrule
  \end{tabularx}
\end{table}

\subsection{Diagrama general}
\label{sec:flujo:diagrama}

El diagrama renderizable se encuentra en
\texttt{diagramas/flujo-productivo.mmd} (mermaid). Para incluirlo en el
PDF se exporta a PNG y se inserta:

\begin{verbatim}
mmdc -i diagramas/flujo-productivo.mmd -o diagramas/flujo-productivo.png
\end{verbatim}

\subsection{Puntos de registro hacia la base de datos (Tópico~7)}
\label{sec:flujo:registro}

Cada etapa con \emph{flecha hacia DB} en el diagrama
\texttt{flujo-productivo.mmd} corresponde a un registro persistente:

\begin{table}[H]
  \centering
  \caption{Puntos de registro a la base de datos unificada.}
  \label{tab:registro}
  \begin{tabularx}{\linewidth}{lXl}
    \toprule
    \textbf{Etapa} & \textbf{Variable registrada} & \textbf{Tabla destino} \\
    \midrule
    1 & Lote, especie, sustrato, masa hidratada                   & \texttt{lotes} \\
    3 & T y t de esterilización, presión registrada               & \texttt{lotes} \\
    4 & Fecha inoculación, jeringa/spawn usado                    & \texttt{lotes} \\
    5 & Detección colonización completa (inspección)              & \texttt{lotes} \\
    6 & Fecha apertura, contenedor destino                        & \texttt{lotes} \\
    7 & Telemetría continua HR, T, CO$_2$, lux (30\,s)            & \texttt{mediciones} \\
    7 & Imagen cada 15\,min                                       & \texttt{imagenes} \\
    8 & Masa fresca cosechada, BE\,\% calculado                   & \texttt{cosechas} \\
    9a/9b & Decisión fresco/seco; masa fresca destinada a secado  & \texttt{cosechas} \\
    9b & Masa seca final, factor conversión, HR residual          & \texttt{cosechas} \\
    \bottomrule
  \end{tabularx}
\end{table}

\subsection{Criterios de aceptación}
\label{sec:flujo:aceptacion}

\begin{itemize}[leftmargin=*,itemsep=0pt]
  \item El diagrama identifica todo punto de registro y lo mapea a una
        tabla del schema (Tópico~7).
  \item Cada transición tiene una condición de paso explícita
        (colonización completa, primordios visibles, sombrero
        aplanado, HR residual $<$ objetivo).
  \item Los tiempos de tabla~\ref{tab:flujo} se actualizan con datos
        propios tras los primeros 3 lotes.
\end{itemize}

% =====================================================================
% §1.6 — Integración de la deshidratadora en el flujo post-cosecha
% =====================================================================
\section{Integración de la deshidratadora en el flujo post-cosecha}
\label{sec:postcosecha}

\subsection*{Resumen de la sección}
El operador cuenta con una deshidratadora doméstica preexistente. Esta
sección queda preparada como \emph{plantilla} para registrar curvas de
secado caracterizadas experimentalmente, referencias bibliográficas de
arranque, criterios de humedad residual objetivo, presentación de
producto seco, y los hooks que conectan el lote a la base de datos
unificada del Tópico~7.

% ---------------------------------------------------------------------
\subsection{Especificaciones de la deshidratadora}
\label{sec:postcosecha:equipo}

% TODO[Operador]:
%   - Marca y modelo exactos.
%   - Potencia eléctrica [W].
%   - Número y área de bandejas [m^2 totales].
%   - Rango térmico operativo declarado por fabricante [°C].
%   - Tipo de flujo (vertical / horizontal).
%   - Sensor interno de temperatura (sí/no, precisión declarada).

\subsection{Parámetros de referencia bibliográfica}
\label{sec:postcosecha:referencias}

% TODO[Fase 1 — relleno]:
% Tabla con valores reportados en literatura, NO experimentales:
%
% | Especie | T secado [°C] | t aprox. [h] | HR residual obj. [%] | Fuente |
% |---------|---------------|--------------|----------------------|--------|
% | Oyster  | 40--50        | 6--10        | <10                  | \cite{} |
% | Lion's mane | 40--55    | 8--12        | <10                  | \cite{} |
% | Shiitake| 50--60        | 8--14        | <12                  | \cite{} |
%
% Estos valores son PUNTO DE PARTIDA. La caracterización propia se
% registra en plantilla-secado.csv y luego se vuelca a la tabla §1.6.4.
%
% Citas mínimas requeridas en referencias.bib:
%   - Chang \& Miles, "Mushrooms: Cultivation, Nutritional Value,
%     Medicinal Effect, and Environmental Impact".
%   - Stamets, "Growing Gourmet and Medicinal Mushrooms".
%   - Algún paper indexado por especie (a buscar).

\subsection{Criterio de humedad residual objetivo}
\label{sec:postcosecha:hr_residual}

% TODO[Fase 1 — relleno]:
% - Justificar rango objetivo <10--12% para conservación estable a
%   temperatura ambiente sin atmósfera modificada.
% - Indicar método de medición:
%     * Indirecto: balanza + masa antes/después.
%     * Directo (opcional, futuro): higrómetro de aw o método AOAC.
% - Validar criterio contra normativa aplicable (NMX, FDA, Codex)
%   — esta validación queda fuera de alcance Fase 1, solo se cita.

\subsection{Plantilla de caracterización (vacía)}
\label{sec:postcosecha:plantilla}

La plantilla operativa vive en
\texttt{plantilla-secado.csv} (raíz del directorio de Fase 1). Se llena
manualmente tras cada lote y, posteriormente, se importa al sistema
del Tópico~7 para construir la curva propia.

% TODO[Fase 1 — relleno]:
% Incluir aquí la misma tabla en LaTeX (vacía) para referencia impresa:
% columnas: Lote, Especie, T [°C], t [h], m_inicial [g], m_final [g],
%           %HR_residual estimada, Observaciones organolépticas.

\subsection{Flujo operativo post-cosecha}
\label{sec:postcosecha:flujo}

% TODO[Fase 1 — relleno]:
% Incluir como figura el diagrama mermaid
% diagramas/flujo-postcosecha.mmd con bifurcación fresco/seco.
% Pasos a representar:
%   - Cosecha → cepillado en seco (NO lavar con agua antes de secar)
%   - Decisión por demanda y antigüedad:
%       fresco → empaque refrigerado → distribución <72 h
%       seco   → deshidratado §1.6 → empaque hermético → distribución
%   - Registro de masa fresca cosechada en lote (Fase 7)

\subsection{Presentación de producto seco}
\label{sec:postcosecha:presentacion}

% TODO[Fase 1 — relleno]:
%   - Formatos: bolsa metalizada sellada al vacío opcional, frasco de
%     vidrio para premium.
%   - Gramajes comerciales sugeridos: 20 g, 50 g, 100 g.
%   - Etiquetado mínimo NO regulatorio (lote, fecha de deshidratado,
%     especie, recomendación de rehidratación).
%   - El detalle regulatorio (NOM-051, registros sanitarios) queda
%     fuera de alcance de la Fase 1.

\subsection{Hooks para la base de datos (Tópico 7)}
\label{sec:postcosecha:hooks}

Para que la curva propia se construya con el tiempo, la tabla
\texttt{cosechas} del schema unificado debe contemplar, por lote:

\begin{itemize}[leftmargin=*]
  \item \texttt{kg\_fresco\_cosechado}
  \item \texttt{kg\_destinado\_a\_secado}
  \item \texttt{kg\_producto\_seco\_final}
  \item \texttt{factor\_conversion} (calculado: seco / destinado)
  \item \texttt{t\_secado\_horas}, \texttt{temp\_secado\_C}
  \item \texttt{hr\_residual\_estimada\_pct}
\end{itemize}

\subsection{Criterios de aceptación}
\label{sec:postcosecha:aceptacion}

\begin{itemize}[leftmargin=*]
  \item Plantilla CSV creada y referenciada.
  \item Al menos 3 referencias bibliográficas (una por especie) en
        \texttt{referencias.bib}.
  \item Diagrama de flujo post-cosecha con bifurcación
        fresco/seco renderizable.
  \item Los campos de la tabla \texttt{cosechas} aparecen como
        requerimiento explícito para el schema del Tópico~7.
\end{itemize}

% =====================================================================
% §1.7 — Preparación de sustrato e inoculación
% =====================================================================
\section{Preparación de sustrato e inoculación}
\label{sec:preparacion}

Subsistema \emph{aguas arriba} del rack: prepara y esteriliza el
sustrato, propaga el micelio en grano (\emph{grain spawn}) y arma las
bolsas inoculadas que se incubarán y trasladarán a los contenedores de
fructificación (\S\ref{sec:contenedores}).

\textbf{Nota.} Esta sección reemplaza el supuesto inicial del brief de
``recibir bolsas ya inoculadas''. El operador realiza el proceso
completo desde sustrato bruto.

% ---------------------------------------------------------------------
\subsection{Esterilización por olla de presión}
\label{sec:preparacion:autoclave}

Se usa una olla de presión doméstica de \paramVolOllaPresion\,L como
autoclave de bajo costo. Parámetros de operación:

\begin{table}[H]
  \centering
  \caption{Ciclos de esterilización en olla de presión.}
  \label{tab:autoclave}
  \begin{tabularx}{\linewidth}{Xcc}
    \toprule
    \textbf{Material} & \textbf{Presión [psi]} & \textbf{Tiempo [min]} \\
    \midrule
    Grano para spawn (frasco 1\,L)               & \paramPresionEster & \paramTiempoEsterGrano \\
    Sustrato en bolsa (paja, pellets, mezclas)   & \paramPresionEster & \paramTiempoEsterSustrato \\
    Material de inoculación (bisturí, jeringas)  & \paramPresionEster & 30 \\
    \bottomrule
  \end{tabularx}
\end{table}

\paragraph{Validación de ciclo.} El tiempo \emph{efectivo} comienza
cuando la válvula libera vapor sostenido; el termómetro de aguja
montado en la tapa permite verificación cruzada. Para sustratos densos
o cargas mayores se incrementa el tiempo en 30\,min por cada nivel
adicional de carga.

\paragraph{Enfriado.} La bolsa o frasco esterilizado debe enfriarse a
$T \leq 30\,\si{\celsius}$ antes de inoculación (típicamente 8--12\,h en
ambiente cerrado). \textbf{No abrir} la bolsa antes de inoculación.

\subsection{Still Air Box (SAB)}
\label{sec:preparacion:sab}

Cámara de aire estático construida con tote hermético de 60--80\,L
adicional, con dos perforaciones circulares (diámetro $\approx$
$\sfrac{1}{2}$ antebrazo) en la pared corta y manga elástica opcional.
Procedimiento de uso:

\begin{enumerate}[leftmargin=*,itemsep=0pt]
  \item Limpieza con peróxido 3\,\% + alcohol isopropílico 70\,\%
        interior y exterior.
  \item Reposo cerrado 5--10\,min para sedimentación de aerosoles.
  \item Introducir bolsa esterilizada, jeringa de spawn, bisturí
        flameado, alcohol en spray.
  \item Inocular minimizando movimientos amplios; sellar la perforación
        del filtro con cinta micropore inmediatamente tras inyectar.
\end{enumerate}

\pendientecurso{Validar dimensiones exactas del SAB y procedimiento
contra recomendaciones del curso. Considerar upgrade a flow hood HEPA
casero si la tasa de contaminación supera 10\,\% en los primeros 5
lotes.}

\subsection{Producción de grain spawn}
\label{sec:preparacion:spawn}

El grain spawn (micelio sobre grano) es el inóculo intermedio entre el
cultivo madre (jeringa de esporas o cultivo en placa) y el sustrato
final. Proceso:

\begin{enumerate}[leftmargin=*,itemsep=0pt]
  \item Enjuagar y remojar grano (quinoa, trigo o mijo) 12--24\,h.
  \item Escurrir y secar superficialmente.
  \item Frasco 1\,L cargado a $\sfrac{2}{3}$; tapa con perforación
        para intercambio gaseoso cubierta con cinta micropore.
  \item Esterilizar en olla \paramVolOllaPresion\,L: \paramTiempoEsterGrano\,min
        a \paramPresionEster\,psi.
  \item Enfriar a $T \leq 30\,\si{\celsius}$ (mínimo 8\,h).
  \item Inocular con jeringa de esporas o transferir trozos de placa
        en SAB.
  \item Incubar en oscuridad a 22--24\,\si{\celsius} hasta colonización
        completa (10--21\,días según especie).
  \item Agitar el frasco una vez a la mitad de la colonización para
        distribuir micelio.
\end{enumerate}

\subsection{Armado de bolsas de sustrato inoculado}
\label{sec:preparacion:bolsas}

\begin{enumerate}[leftmargin=*,itemsep=0pt]
  \item Hidratar sustrato a campo (HR $\approx 65$\,\%) y agregar
        suplemento (yeso, salvado) según especie.
  \item Embolsar 1{,}5--2{,}5\,kg por bolsa autoclavable con parche
        filtro 0{,}2\,\si{\micro\meter}.
  \item Esterilizar a \paramPresionEster\,psi durante
        \paramTiempoEsterSustrato\,min.
  \item Enfriar.
  \item Inocular en SAB con $\approx 5$--$10$\,\% del peso húmedo en
        grain spawn.
  \item Cerrar bolsa con sello impulso o liga doblada.
  \item Etiquetar (lote, especie, fecha) y trasladar a incubación.
\end{enumerate}

\subsection{Control de contaminación}
\label{sec:preparacion:contam}

Registro obligatorio por lote: número de bolsas inoculadas, número de
bolsas con contaminación detectada (Trichoderma verde, mohos negros,
bacterial sour), fecha de detección. Métricas objetivo:

\begin{itemize}[leftmargin=*,itemsep=0pt]
  \item Tasa de contaminación $<$10\,\% en los primeros 5 lotes.
  \item Tasa $<$5\,\% en régimen establecido.
  \item Una tasa $>$25\,\% indica revisar: tiempo de esterilización,
        sello del filtro de la bolsa, técnica en SAB o calidad del
        cultivo madre.
\end{itemize}

\subsection{Hooks para la base de datos (Tópico~7)}
\label{sec:preparacion:hooks}

Campos adicionales en la tabla \texttt{lotes}:

\begin{itemize}[leftmargin=*,itemsep=0pt]
  \item \texttt{sustrato\_tipo}, \texttt{sustrato\_masa\_hidratada\_kg}
  \item \texttt{esteriliz\_temp\_C}, \texttt{esteriliz\_tiempo\_min},
        \texttt{esteriliz\_presion\_psi}
  \item \texttt{spawn\_origen} (jeringa lote / placa propia)
  \item \texttt{spawn\_porcentaje} (\% sobre masa húmeda)
  \item \texttt{contaminacion\_detectada\_bool},
        \texttt{contaminacion\_tipo} (enum)
  \item \texttt{fecha\_inoculacion}, \texttt{fecha\_colonizacion\_completa}
\end{itemize}

\subsection{Criterios de aceptación}
\label{sec:preparacion:aceptacion}

\begin{itemize}[leftmargin=*,itemsep=0pt]
  \item Olla de presión instalada, ciclos verificados con termómetro
        independiente.
  \item SAB armado y sanitizado.
  \item Primer lote de grain spawn colonizado a $>$90\,\%.
  \item Tasa de contaminación documentada para cada lote.
  \item Campos en \texttt{lotes} reflejados en el schema del Tópico~7.
\end{itemize}

% ---------------------------------------------------------------------

\printbibliography

\end{document}
